\section{Dependence, sums and generating functions}
\subsection*{Part I. Joint distributions and conditional expectation}
\begin{problembook}
We independently roll two red and one blue standard dice. Let $X$ be the number
of sixes on the red dice, and $Y$ the number of twos on all three dice. Find the
joint distribution and $\Pbb(X>0\mid Y=1)$. Are $X$ and $Y$ independent?
Can you answer the last question without calculating the entire joint distribution?
\end{problembook}
\begin{problembook}
A fair coin is tossed independently four times. Let $X$ be the number of heads
in the first three tosses, and $Y$ the number in the last two.
\begin{enumeratea}
  \item Find the joint and conditional distributions, $\Pbb(X=Y)$ and $\Pbb(X>1\mid Y=1)$.
  \item Find $\Ebb(X\mid Y)$ and $\Ebb(Y\mid X)$ and their distributions.
  Verify that $\Ebb[\Ebb(X\mid Y)]=\Ebb X$.
  \item Find $\Cov(X,Y)$ in two ways: from the joint distribution, and by writing
  $X$ and $Y$ as sums of indicators.
\end{enumeratea}
\end{problembook}
\begin{problembook}
Let $X$ be uniformly distributed on $\{-1,0,1\}$ and $Y=X^2$.
\begin{enumeratea}
  \item Show that $\Cov(X,Y)=0$, but $X$ and $Y$ are dependent.
  \item Now let $X$ and $Y$ take values in $\{0,1\}$. Prove that $\Cov(X,Y)=0$
  is equivalent to independence. Why does this not contradict the preceding example?
  \item For random variables $X,Y$ with finite positive variances, prove that
  $|\Cor(X,Y)|\le1$. When does equality hold?
\end{enumeratea}
\end{problembook}
\begin{problempencil}
One box contains 4 white and 6 black balls, and another contains 7 white and
3 black balls. From each box we draw one ball uniformly, independently of the
other draw. We add one white ball to the two drawn balls and place all three in a third box.
\begin{enumeratea}
  \item What is the probability that a uniformly chosen ball from the third box is white?
  \item Draw from the third box with replacement until the first white ball.
  For a fixed box composition, draws are independent and uniform. Let $X$ be the
  number of draws. Find $\Ebb X$. Does $X$ have a geometric distribution?
  \item Find $\Pbb(X>k)$ for every integer $k\ge0$, and prove that
  $\Pbb(X>k+1\mid X>k)$ increases with $k$. What do we learn about the box's
  composition after a sequence of black balls?
\end{enumeratea}
\end{problempencil}
\subsection*{Additional problems}
\begin{problembook}
Choose three numbers uniformly without replacement from $\{1,2,3,4,5\}$.
Let $X$ be the middle number in order of size, and $Y$ the smallest.
\begin{enumeratea}
  \item Find the joint and marginal distributions. Are $X,Y$ independent?
  \item Find $\Cov(X,Y)$ and $\Cor(X,Y)$.
  \item Find the distribution, expectation and variance of $X-2Y$.
\end{enumeratea}
\end{problembook}
\begin{problempencil}
Let $(X_1,\ldots,X_n)$ be a uniformly random permutation of $\{1,\ldots,n\}$,
where $n\ge2$, and let $S=X_1+\cdots+X_n$.
\begin{enumeratea}
  \item Find $\Ebb S$, $\Var S$, $\Ebb X_i$ and $\Var X_i$.
  \item Prove the variance-of-a-sum formula and use it to find $\Cov(X_i,X_j)$
  for $i\ne j$. Explain the sign of the covariance.
  \item For $1\le k\le n$, find $\Var(X_1+\cdots+X_k)$. Compare with the variance
  of $k$ independent draws with replacement.
\end{enumeratea}
\end{problempencil}
\begin{problempencil}
On each face of a blank die, independently write a uniformly chosen number from
1 to 6. Then roll the resulting die $n\ge2$ times. The faces rolled are independent,
uniformly selected and independent of the labels. Let $T_k$ be the number seen on roll $k$.
\begin{enumeratea}
  \item What is the distribution of $T_1$? Find $\Cor(T_1,T_2)$. Are $T_1,T_2$ independent?
  \item If $X=T_1+\cdots+T_{n-1}$ and $Y=T_2+\cdots+T_n$, find $\Cor(X,Y)$.
  \item How do the two correlations change if the labels are independent and uniform
  on $\{199,\ldots,999\}$? What if the labels are independent with an arbitrary
  common distribution having finite positive variance?
\end{enumeratea}
\end{problempencil}
\begin{problem}[P\'olya's urn]
An urn contains one red and one black ball. At each step, choose a ball uniformly,
replace it and add another ball of the same colour. Let $X_n$ indicate a red ball
on draw $n$, and let $W_n=X_1+\cdots+X_n$. Find $\Pbb(X_n=1)$, $\Cov(X_i,X_j)$
for $i\ne j$, and the distribution of $W_4$. What distribution do you conjecture for $W_n$?
\end{problem}
\subsection*{Challenges}
\begin{problemstar}[Pairwise independence]
Let $q$ be prime, and let $A,B$ be independent and uniform on $\{0,\ldots,q-1\}$.
All subsequent calculations are modulo $q$. For $t\in\{0,\ldots,q-1\}$, define
\[ X_t=A+tB\pmod q,\qquad X_\infty=B. \]
\begin{enumeratea}
  \item Prove that $X_0,\ldots,X_{q-1},X_\infty$ are uniform and pairwise
  independent, but not mutually independent. What do the values of two distinct coordinates determine?
  \item Let $N$ be the number of zero coordinates in this family.
  Find its distribution, expectation and variance.
  \item Calculate the variance using indicators as well. Is pairwise independence
  enough for variances to add?
\end{enumeratea}
\end{problemstar}
\begin{problemstar}
For random variables $X,Y$ with finite second moments and a random variable $Z$, prove that
\[
\Cov(X,Y)=\Ebb\!\left[\Cov(X,Y\mid Z)\right]
+\Cov\!\left(\Ebb[X\mid Z],\Ebb[Y\mid Z]\right).
\]
What formula results when $Y=X$? Now let $Z\sim\Ber(1/2)$, let $U$ be uniform
on $\{-1,1\}$, let $Z,U$ be independent, and define
\[ X=3Z+U,\qquad Y=3Z-U. \]
\begin{enumeratea}
  \item Find the conditional covariance and correlation for fixed $Z$.
  \item Find $\Cov(X,Y)$ and $\Cor(X,Y)$.
  \item How do you explain the difference in signs? Relate your answer to the two terms in the formula.
\end{enumeratea}
\end{problemstar}
\subsection*{Part II. Generating functions and random sums}
\begin{problempencil}[Test 1, CS 2024]
For a nonnegative integer-valued random variable $X$, write $g_X(s)=\Ebb[s^X]$, $0\le s\le1$.
\begin{enumeratea}
  \item Let $X$ be uniform on $\{1,\ldots,6\}$. Find $g_X$ and express
  $\Ebb X$ and $\Var X$ in terms of its derivatives.
  \item If $X,Y$ are independent nonnegative integer-valued random variables,
  prove that $g_{X+Y}=g_Xg_Y$.
  \item[$(\mathrm{c})^\star$] Is the converse true?
\end{enumeratea}
\end{problempencil}
\begin{problembook}
Let $X,Y$ be independent. Use generating functions to find the distribution of
$X+Y$ when $X\sim\Bin(n,p)$ and $Y\sim\Bin(m,p)$, and then when
$X\sim\Poi(\lambda)$ and $Y\sim\Poi(\mu)$. Similarly find the distribution of
$G_1+\cdots+G_r$, where $G_1,\ldots,G_r$ are independent and $G_i-1\sim\Geom(p)$,
so each $G_i$ counts the trials up to the first success. Why does equality of
generating functions determine the distribution uniquely? Can the sum of independent
$\Bin(n,p)$ and $\Bin(m,q)$ variables be binomial when $n,m\ge1$, $0<p,q<1$ and $p\ne q$?
\end{problembook}
\begin{problempencil}
The number of visitors to a stadium is $N\sim\Poi(\lambda)$. Each visitor selects
one of ten entrances uniformly, independently of everyone else and of $N$.
\begin{enumeratea}
  \item Find the distribution, expectation and variance of the number using the first entrance.
  \item More generally, let $N,X_1,X_2,\ldots$ be independent nonnegative
  integer-valued random variables, with $(X_i)$ identically distributed.
  For $S=X_1+\cdots+X_N$, setting $S=0$ when $N=0$, prove that
  \[ g_S(s)=g_N\bigl(g_{X_1}(s)\bigr). \]
  If $N,X_1$ have finite second moments, find $\Ebb S$ and $\Var S$.
  \item Apply the result to the stadium problem. Also obtain the distribution
  in (a) using the law of total probability.
\end{enumeratea}
\end{problempencil}
\begin{problembook}
$A$ and $B$ play independent games. In each game, $A$ wins with probability $2/3$
and $B$ with probability $1/3$. Play stops when someone wins two consecutive games.
Let $T$ be the number of games played. Find $\Ebb T$ and the probability that
$A$ wins the contest. What information about the past suffices to describe future
play? Also find the generating function of $T$.
\end{problembook}
\subsection*{Additional problems}
\begin{problempencil}
Two teams each take one penalty in each round. All attempts are independent,
with scoring probabilities $0.75$ for the first team and $0.80$ for the second.
All five regular rounds are played. If tied afterwards, play continues until,
in one round, one team scores and the other does not.
\begin{enumeratea}
  \item What is the probability that in each of the first five rounds both teams
  have the same outcome for their penalty?
  \item What is the probability that play continues after the fifth round? After the tenth?
  \item Find the expected total number of penalties taken.
\end{enumeratea}
\end{problempencil}
\begin{problem}
Lottery tickets are numbered from 000000 to 999999. Find the probability that a uniformly chosen ticket:
\begin{enumeratea}
  \item has digit sum 21;
  \item has equal sums of its first and last three digits;
  \item has a sum of its first three digits that exceeds that of its last three by 2.
\end{enumeratea}
It suffices to express the answers as coefficients of a polynomial.
\end{problem}
\begin{problempencil}
We telephone $N\sim\Poi(\lambda)$ people with an offer. Person $i$ would accept
on the $K_i$th call, where $K_1,K_2,\ldots$ are independent, identically distributed
and independent of $N$, with $\Pbb(K_i=n)=f(n)$ for $n\ge1$ and
$\sum_{n\ge1}f(n)=1$. Let $X_n$ be the number of people who would accept on call $n$.
\begin{enumeratea}
  \item Prove that $(X_n)_{n\ge1}$ are independent and $X_n\sim\Poi(\lambda f(n))$.
  \item We stop at the first round with no new acceptance. If
  $T=\inf\{n\ge1:X_n=0\}$ and $S=X_1+\cdots+X_T$, prove that
  \[ \Ebb S=\lambda\Ebb[F(T)],\qquad F(k)=\sum_{j=1}^k f(j). \]
  Why can we not directly apply the expectation formula for a random sum whose
  number of summands is independent of the summands?
\end{enumeratea}
\end{problempencil}
\begin{problempencil}
Two standard dice are rolled independently. Can we label their faces with other
positive integers so that independent rolls have the same sum distribution?
Repeated labels are allowed. Find a construction different from the two standard dice.
\end{problempencil}
\subsection*{Challenges}
\begin{problemstar}\label{prob:week4-penney}
We toss a fair coin independently, recording outcomes as $H,T$. For a word
$w=w_1w_2w_3$, let $\tau_w$ be the toss number when it first appears as three
consecutive outcomes. Consider $a=HHT$ and $b=THH$.
\begin{enumeratea}
  \item Choose states retaining enough information about the appearance of the
  two words, and draw the possible transitions.
  \item Find $\Pbb(\tau_a<\tau_b)$ and $\Ebb(\tau_a\wedge\tau_b)$.
  \item Prove that $\tau_a\wedge\tau_b<\infty$ almost surely, without using the calculated expectation.
\end{enumeratea}
\end{problemstar}
\begin{problemstar}[Branching process]
A population starts with one individual. Each individual has 0 children with
probability $1-p$ and 2 children with probability $p$, where $0<p<1$, independently
of all other individuals in all generations. Let $Z_n$ be the number of individuals
in generation $n$, with $Z_0=1$, and let
\[ f(s)=1-p+ps^2,\qquad T=\inf\{n\ge0:Z_n=0\}. \]
\begin{enumeratea}
  \item Prove that $g_{Z_n}=f^{\circ n}$.
  \item Show that $q_n=\Pbb(T\le n)=f^{\circ n}(0)$ and that $(q_n)$ is increasing.
  \item Prove that $q=\Pbb(T<\infty)$ is the smallest solution of $f(q)=q$ in
  $[0,1]$, and find it as a function of $p$.
  \item Find $\Ebb Z_n$. When $p=1/2$, it is constant, yet extinction occurs almost surely.
  How do you explain this?
\end{enumeratea}
\end{problemstar}
\subsection*{Independent research}
\begin{researchtask}[Penney's game]
Two players each choose a different word from $\{H,T\}^3$. We then toss a fair
coin independently until one of the words appears. The player whose word appears first wins.
\begin{enumeratea}
  \item Generalise the model in Problem~\ref{prob:week4-penney} and write a program
  that calculates exactly, for each of the 56 ordered pairs $(a,b)$,
  \[ p(a,b)=\Pbb(\tau_a<\tau_b),\qquad e(a,b)=\Ebb(\tau_a\wedge\tau_b). \]
  Use rational arithmetic and check at least one pair by hand.
  \item For each first word $a$, find the best responses $b$ for the second player.
  Does the player choosing second have an advantage? Formulate a general rule
  for the best response and prove it.
  \item Draw a directed graph with edge $a\to b$ when $b$ is a best response to
  $a$. Find a cycle and explain what it shows about comparisons between words.
\end{enumeratea}
\end{researchtask}
\noindent\textbf{Lean.} Formalise in Lean 4 with Mathlib and prove that, for
finite sets $S,T$ and functions $a:S\to\mathbb R$ and $b:T\to\mathbb R$,
\[
\sum_{s\in S}\sum_{t\in T}a(s)b(t)
=\left(\sum_{s\in S}a(s)\right)\left(\sum_{t\in T}b(t)\right).
\]
Explain how this identity is used in proving $g_{X+Y}=g_Xg_Y$ for independent
nonnegative integer-valued random variables $X,Y$ with finite ranges.
